\documentclass[12pt]{article}
\usepackage{amsmath, amssymb, amsthm}
\usepackage{mathtools}
\usepackage{graphicx}
\usepackage{float}
\usepackage{hyperref}
\usepackage{xcolor}
\usepackage{listings}
\usepackage{geometry}
\usepackage{algorithm}
\usepackage{algpseudocode}
\usepackage{tikz}
\usepackage{longtable}
\usepackage{circuitikz}
\usepackage{comment}
\usepackage{MnSymbol}
\usepackage{physics}
\usepackage{booktabs}
\usepackage{subfig} % For \subfloat
\usepackage[section]{placeins}

\begin{document}

\begin{titlepage}
    \centering
    \vspace*{1.5cm}
    
    {\Huge\bfseries Group Quiz 2}\\[0.5cm]
    {\Large\bfseries Differential Equations and Transform Techniques (EE1060)}\\[1cm]
    {\Large\bfseries Detailed Analysis of Convolution for Continous Time Kernel}\\[1.5cm]
    
    \rule{0.8\textwidth}{0.4pt}\\[0.5cm]
    {\Large\bfseries Submitted By}\\[0.5cm]
    \rule{0.8\textwidth}{0.4pt}\\[0.5cm]
    
    \large
    Krishna Patil \hfill EE24BTECH11036\\[0.2cm]
    Agamjot Singh \hfill EE24BTECH11002\\[0.2cm]
    Kollaru Suraj \hfill EE24BTECH11033\\[0.2cm]
    Dwanith M Doddahundi \hfill EE24BTECH11016\\[0.2cm]
    Swaraj Sharma Durgi \hfill EE24BTECH11018\\[0.2cm]
    Nevanath Vasu \hfill EE24BTECH11046\\[1.5cm]
    
\end{titlepage}


\tableofcontents
\newpage

\section{Problem Statement}
\begin{enumerate}
\item Compute the convolution of a given signal $f(t)$ with a rectangular kernel $h(t)$, analytically. The rectangular kernel is defined as:

\begin{equation}
h(t) =
\begin{cases}
1, & \text{for } -T \leq t \leq T \\
0, & \text{otherwise}
\end{cases}
\end{equation}

Derive the convolution expression $y(t) = (f * h)(t)$ in terms of known functions, and analyze the system's behavior for various values of the kernel duration $T$ and the input signal $f(t)$. Additionally, investigate the following scenarios:

\begin{enumerate}
\item Modify the kernel to only consider the part of the kernel for $t > 0$. How does this affect the convolution result?

\item Shift the kernel by a time $\tau_0$. Analyze how the shift impacts the convolution output and discuss the significance of this shift in the context of time-delayed systems.
\end{enumerate}

You are free to choose the form of the input signal $f(t)$, but make sure it is well-defined and appropriate for convolution. A step or sinusoidal signal would work well for this analysis.
\end{enumerate}

\input{step.tex}
\input{exp.tex}
\input{sin_1.tex}
\input{sin_2.tex}
\input{sin_3.tex}
\input{smoothening.tex}

\end{document}
